\section{Conclusion}
\label{sec:conclusion}
The transport coefficients that admit a single-particle Green--Kubo expansion are coarse linear functionals of one object, the atom-wise integrated correlation matrix $C$, and its spectrum can be read against a random-matrix null once the sample count of the estimator is set by the integration lag rather than by the decorrelation time alone.
With that null in place, two results stand out.

First, bonding has a precise spectral signature.
Because $C$ is the long-time displacement covariance, every independent bounded relative coordinate sends one eigenvalue below the noise bulk, with an eigenvector on that coordinate.
In the Lennard-Jones liquid this holds bond for bond --- $8$ and $27$ sub-bulk eigenvalues for $8$ and $27$ bonds in every seed, with eigenvectors that recover $86$--$94\%$ of the bond-constraint space --- provided the integration lag exceeds the time a free pair needs to diffuse over the bond length.
At shorter lags the same information is present in the pair-resolved elements of $C$ but does not separate from the bulk; the criterion $r_{0}^{2}/(6D\tau_{\max}) < (1-\sqrt{q})^{2}$ predicts the measured bonded mode to within about $15\%$ and makes the lag a resolution parameter.

Second, the entropy of $C$ measured against the null distinguishes states of matter.
Along a Lennard-Jones isobar the entropy deficit separates crystal, disordered solid and liquid by an order of magnitude and decays smoothly towards zero within the liquid.
On the temperature grid used here it does not resolve the transitions more sharply than the density, and in the solids its structure reflects vibrational correlation within the integration window rather than transport.

The same two signatures carry over to molecular and ionic liquids.
In water, the rigid geometry of each molecule places exactly two eigenvalues per molecule far below the molecular bulk, with eigenvectors that span the intramolecular coordinates to better than $97\%$, and the entropy deficit of the oxygen motion separates ice from liquid and rises continuously as the liquid is supercooled.
In molten NaCl, the Nernst--Einstein deviation read from the species blocks of $C$ is consistent with the charge-current Green--Kubo integral, as it must be by construction, but its split into self and distinct parts converges slowly with the lag, whereas the conductivity itself is robust; the entropy deficit of the liquid reflects short-lived cation--anion correlation and vanishes at lags of a few picoseconds.

Two extensions follow directly.
The decomposition applies unchanged to shear viscosity and thermal conductivity once the pair contribution to their fluxes is partitioned atom by atom, and because a bond registers only if it outlives the lag, the lag dependence of the sub-bulk spectrum may give access to bond-lifetime distributions in associating liquids; all bonds studied here are permanent, so this remains to be tested.
